\paragraph{A same-code reference.}
The controlled panel contains nine distinct members; its reference
contains nine byte-identical \bare{} slots. Each member or slot receives
three executions per task. This design preserves population size and
sampling opportunities while changing program identity.

\paragraph{Learning a choice on two repeats, testing it on the third.}
\label{sec:calibration-test}
Each fold uses two discovery repeats to select a member for every task
and a single best fixed member. The third repeat measures their accuracy
difference. Averaging the three folds gives $G_{\rm real}$ and
$G_{\rm clone}$. We compare paired task-level gains on the common
complete tasks through
\begin{equation}
D=G_{\rm real}-G_{\rm clone}.
\label{eq:clone-difference}
\end{equation}
A positive difference means that choices learned from past outcomes
transfer better for distinct programs than for identical ones.
Both the task-wise choice and the fixed comparator are determined
entirely by the discovery repeats.

\paragraph{Analysis set and repeatability.}
All primary estimates use the same $\RcvPairedN$ tasks, with
$\RcvCells$ scored cells per arm. The original incomplete-task fraction
is $\RcvOriginalExcludedPct\%$, below the acquisition gate's $10\%$
ceiling. The acquisition protocol preceded collection; the repeat
diagnostic and common-complete-set analysis were developed afterward.
The recovery snapshot preserves the programs, solver, tasks, and
three-repeat design (Appendix~\ref{app:wp1r}).

We also measure residual score-pattern correlation across repeats
after removing additive member and task effects. This exploratory
measure describes repeatability, while $G$ describes held-out selection
gain. Known-truth simulations characterize the diagnostic's sensitivity.
Appendices~\ref{app:repeatability} and~\ref{app:statistics} specify these
calculations, tie handling, uncertainty estimates, and decision rules.
